% content.tex  —  leanblueprint for Peter–Weyl / non-abelian Fourier analysis on
% finite groups.  Following T. Tao,
% https://terrytao.wordpress.com/2011/01/23/the-peter-weyl-theorem-and-non-abelian-fourier-analysis-on-compact-groups/
% specialised to the finite-group case throughout.
%
% Design principles for autoformalisation:
%   * atomic lemmas: each step that can be a self-contained Lean declaration
%     is its own \begin{lemma}…\end{lemma} with a small, well-typed signature.
%   * every \lean{...} reference is a real Mathlib name unless tagged [NEW].
%   * Mathlib availability key in comments:
%       [ML]  = already in Mathlib; \mathlibok and \leanok appropriate.
%       [~ML] = close analogue exists in Mathlib; minor adaptation needed.
%       [NEW] = not yet in Mathlib; needs local proof.
%   * we use Mathlib's bundled API throughout:
%       - representations: Representation k G V (= G →* (V →ₗ[k] V)) and the
%         category Rep k G of bundled representations; finite-dimensional
%         versions live in FDRep k G.
%       - irreducibility: CategoryTheory.Simple V for V : FDRep k G; equivalent
%         to IsSimpleModule (MonoidAlgebra k G) ρ.asModule.
%       - intertwining maps: morphisms V ⟶ W in Rep k G / FDRep k G.
%       - Maschke / semisimplicity:
%         MonoidAlgebra.Submodule.instIsSemisimpleModule
%         MonoidAlgebra.Submodule.exists_isCompl
%         LinearMap.equivariantProjection / sumOfConjugatesEquivariant.
%       - characters: FDRep.character, FDRep.char_orthonormal,
%         FDRep.average_char_eq_finrank_invariants,
%         Representation.card_inv_mul_sum_char_mul_char_eq_finrank.
%
% Throughout this blueprint we fix:
%   * a finite group G (Mathlib: [Group G] [Fintype G] or [Finite G]),
%   * an algebraically closed field k (Mathlib: [Field k] [IsAlgClosed k])
%     in which the order of G is invertible (Mathlib:
%     [Invertible (Fintype.card G : k)] or equivalently [NeZero (Nat.card G : k)]).
%   The classical case k = ℂ satisfies all of these.

\chapter{Background: Finite Groups, Representations, Intertwiners}
\label{chap:background}

\section{Hypotheses and Notation}

\begin{definition}[Standing hypotheses]
  \label{def:standing-hypotheses}
  \mathlibok
  % [ML] Group, Fintype, Field, IsAlgClosed, Invertible all in Mathlib core.
  Throughout the blueprint the symbol $G$ denotes a finite group: in Lean,
  \begin{center}
    \texttt{[Group G] [Fintype G]}.
  \end{center}
  The symbol $k$ denotes an algebraically closed field whose characteristic does
  not divide $|G|$:
  \begin{center}
    \texttt{[Field k] [IsAlgClosed k] [Invertible ((Fintype.card G : k))]}.
  \end{center}
  We write $|G|$ for \texttt{Fintype.card G} and view it as an element of $k$
  via the canonical map $\mathbb Z \to k$.
\end{definition}

\section{Representations and Their Bundlings}

\begin{definition}[$k$-linear representation]
  \label{def:representation}
  \lean{Representation}
  \mathlibok
  % [ML] Mathlib.RepresentationTheory.Basic:
  %   Representation k G V := G →* (V →ₗ[k] V)  as an abbreviation.
  \uses{def:standing-hypotheses}
  A \emph{$k$-linear representation} of $G$ on a $k$-module $V$ is a monoid
  homomorphism $\rho : G \to (V \to_{\!\ell[k]} V)^{\!\times}$ from $G$ into the
  multiplicative monoid of $k$-linear endomorphisms of $V$. In Mathlib this is
  the abbreviation \texttt{Representation k G V := G →* (V →ₗ[k] V)}.
\end{definition}

\begin{definition}[$\rho.\mathrm{asModule}$ : MonoidAlgebra-module structure]
  \label{def:as-module}
  \lean{Representation.asModule}
  \mathlibok
  % [ML] Mathlib.RepresentationTheory.Basic:
  %   ρ.asModule : Type, with [Module (MonoidAlgebra k G) ρ.asModule].
  %   This is V with the k[G]-action (∑_g a_g g) • v = ∑_g a_g · ρ(g) v.
  \uses{def:representation}
  Given $\rho : \mathrm{Representation}\ k\ G\ V$, the type
  \texttt{ρ.asModule} carries a canonical \texttt{MonoidAlgebra k G}-module
  structure that extends the $k$-action on $V$. Conversely, any
  $k[G]$-module gives rise to a representation. This identification (Mathlib
  records it as the equivalence \texttt{Rep.equivalenceModuleMonoidAlgebra}) is
  the bridge between representation-theoretic and module-theoretic statements.
\end{definition}

\begin{definition}[Bundled categories \texttt{Rep k G} and \texttt{FDRep k G}]
  \label{def:rep-categories}
  \lean{Rep, FDRep}
  \mathlibok
  % [ML] Mathlib.RepresentationTheory.Rep, Mathlib.RepresentationTheory.FDRep.
  %   Rep k G   : k-linear abelian symmetric monoidal category of all reps.
  %   FDRep k G : full subcategory of finite-dimensional reps.
  \uses{def:representation}
  \texttt{Rep k G} is the category of $k$-linear $G$-representations:
  objects are pairs $(V, \rho)$ as above, morphisms are $G$-equivariant
  $k$-linear maps. \texttt{FDRep k G} is the full subcategory of finite-dimensional
  representations. For $V : \texttt{FDRep k G}$, the underlying vector space is
  \texttt{V.V} (or by coercion just \texttt{V}) and the action is
  \texttt{V.ρ : G →* (V →ₗ[k] V)}.
\end{definition}

\begin{definition}[Intertwining map = morphism in \texttt{FDRep k G}]
  \label{def:intertwiner}
  \lean{FDRep}
  \mathlibok
  % [ML] In Mathlib, an intertwining map between V W : FDRep k G is just
  %   a morphism V ⟶ W of FDRep k G. Concretely Action.Hom : the underlying
  %   k-linear map together with the equivariance square.
  %   The k-vector space structure on (V ⟶ W) is FDRep.instAddCommGroup +
  %   FDRep.instModule, and finite-dimensionality is FDRep.instFiniteDimensional.
  \uses{def:rep-categories}
  An \emph{intertwining map} (or \emph{$G$-map}) from $V$ to $W$ in
  \texttt{FDRep k G} is a morphism $V \longrightarrow W$, i.e.\ a $k$-linear
  map $T : V \to W$ such that $T \circ \rho_V(g) = \rho_W(g) \circ T$ for all
  $g \in G$. The hom space $V \to W$ is finite-dimensional over $k$.
\end{definition}

\begin{definition}[Isomorphism of representations]
  \label{def:rep-iso}
  \lean{CategoryTheory.Iso}
  \mathlibok
  % [ML] CategoryTheory.Iso V W in Rep k G / FDRep k G is the standard
  %   categorical isomorphism: a pair of mutually-inverse intertwiners.
  \uses{def:intertwiner}
  Two representations $V, W : \texttt{FDRep k G}$ are \emph{isomorphic},
  $V \cong W$, if there is an isomorphism $V \cong W$ in the category
  \texttt{FDRep k G}: equivalently, a $k$-linear isomorphism whose underlying
  map is $G$-equivariant.
\end{definition}

\begin{definition}[Subrepresentation, invariants]
  \label{def:subrep}
  \lean{Representation.invariants}
  \mathlibok
  % [ML] Submodules of ρ.asModule (over MonoidAlgebra k G) are exactly
  %   G-invariant k-subspaces.  The fixed-point subspace is
  %   Representation.invariants ρ : Submodule k V.
  \uses{def:as-module}
  A \emph{subrepresentation} of $\rho$ on $V$ is an
  \texttt{(MonoidAlgebra k G)}-submodule of \texttt{ρ.asModule}; equivalently,
  a $k$-subspace $W \subseteq V$ with $\rho(g)(W) \subseteq W$ for every
  $g \in G$. The \emph{$G$-invariants} are
  $V^G := \{v \in V : \rho(g)v = v\ \forall g\}$, a $k$-subspace
  \texttt{Representation.invariants ρ}.
\end{definition}

\begin{definition}[Direct sum of representations]
  \label{def:direct-sum}
  \mathlibok
  % [ML] FDRep k G has biproducts (categorical (co)products coincide because
  %   the category is abelian k-linear). Concretely V ⊞ W with diagonal action.
  \uses{def:rep-categories}
  The category \texttt{FDRep k G} has biproducts: for $V, W : \texttt{FDRep k G}$
  there is $V \oplus W : \texttt{FDRep k G}$ (notation \texttt{V ⊞ W}) on the
  underlying vector space $V \oplus W$ with the diagonal action.
\end{definition}

\begin{definition}[Irreducible representation]
  \label{def:simple}
  \lean{CategoryTheory.Simple}
  \mathlibok
  % [ML] CategoryTheory.Simple V for V : FDRep k G means V ≠ 0 and every
  %   monic into V is iso or zero.  In our setting this is equivalent to:
  %   V ≠ 0 has no proper non-zero G-invariant subspace.
  %   Equivalently, IsSimpleModule (MonoidAlgebra k G) (V.ρ.asModule).
  \uses{def:rep-categories, def:subrep}
  A finite-dimensional representation $V : \texttt{FDRep k G}$ is
  \emph{irreducible} (or \emph{simple}) if
  \texttt{CategoryTheory.Simple V}: $V$ is non-zero and every subrepresentation
  is either $0$ or $V$ itself. In Lean this is provided as a typeclass
  \texttt{[CategoryTheory.Simple V]}.
\end{definition}

\begin{definition}[Inner product on ${k^G}$]
  \label{def:inner-prod}
  \lean{ClassFunction.innerProduct}
  % [NEW] Not packaged in Mathlib as a standalone bilinear form, but the
  %       defining sum appears explicitly in
  %       Representation.card_inv_mul_sum_char_mul_char_eq_finrank.
  \uses{def:standing-hypotheses}
  For $f, h : G \to k$ define
  \[
    \langle f, h\rangle_G \;:=\; \frac{1}{|G|}\sum_{g \in G} f(g)\, h(g^{-1}).
  \]
  When $h = \overline{h'}$ is replaced by complex conjugation in the unitary
  case, this recovers the usual Hermitian inner product. We use the algebraic
  form throughout, which agrees with Mathlib's character formula
  \texttt{Representation.card\_inv\_mul\_sum\_char\_mul\_char\_eq\_finrank}.
\end{definition}


\chapter{Maschke's Theorem and Semisimplicity}
\label{chap:maschke}

The strategy is the standard ``averaging trick'': start from any $k$-linear
retraction of a $G$-equivariant inclusion, then average over $G$ to get a
$G$-equivariant retraction. We follow the Mathlib formulation in
\texttt{Mathlib.RepresentationTheory.Maschke}.

\section{The Conjugate of a Linear Map}

\begin{definition}[Conjugate of a linear map]
  \label{def:conjugate}
  \lean{LinearMap.conjugate}
  \mathlibok
  % [ML] Mathlib.RepresentationTheory.Maschke:
  %   LinearMap.conjugate : (W →ₗ[k] V) → G → (W →ₗ[k] V)
  %   π.conjugate g v = single g⁻¹ 1 • π (single g 1 • v).
  \uses{def:as-module}
  Let $V, W$ be $k[G]$-modules and $\pi : W \to V$ a $k$-linear map. For each
  $g \in G$ define the \emph{conjugate}
  \[
    (\pi^g)(v) \;:=\; g^{-1} \cdot \pi(g \cdot v),
  \]
  i.e.\ \texttt{π.conjugate g v = (single g⁻¹ 1) \(\bullet\) π((single g 1) \(\bullet\) v)}.
\end{definition}

\begin{lemma}[Conjugate fixes equivariant inclusions]
  \label{lem:conjugate-of-i}
  \lean{LinearMap.conjugate_i}
  \mathlibok
  % [ML] Mathlib.RepresentationTheory.Maschke: LinearMap.conjugate_i.
  \uses{def:conjugate}
  If $i : V \to W$ is $k[G]$-linear and $\pi \circ i = \mathrm{id}_V$, then
  for every $g \in G$ also $\pi^g \circ i = \mathrm{id}_V$ (as $k$-linear maps).
\end{lemma}

\begin{proof}
  \leanok
  \uses{def:conjugate}
  By unfolding: $\pi^g(i(v)) = g^{-1}\cdot\pi(g\cdot i(v)) = g^{-1}\cdot\pi(i(g\cdot v))
  = g^{-1}\cdot(g\cdot v) = v$, where we used $k[G]$-linearity of $i$ and the
  retraction hypothesis. This is \texttt{LinearMap.conjugate\_i} in Mathlib.
\end{proof}

\section{The Averaged (Equivariant) Projection}

\begin{definition}[Sum-of-conjugates]
  \label{def:sum-of-conjugates}
  \lean{LinearMap.sumOfConjugatesEquivariant}
  \mathlibok
  % [ML] Mathlib.RepresentationTheory.Maschke:
  %   LinearMap.sumOfConjugatesEquivariant G π = ∑ g, π.conjugate g.
  \uses{def:conjugate}
  Define $\Sigma\pi := \sum_{g \in G} \pi^g$.  The Mathlib name is
  \texttt{LinearMap.sumOfConjugatesEquivariant G π}, and the corresponding
  pointwise formula is \texttt{sumOfConjugatesEquivariant\_apply}.
\end{definition}

\begin{definition}[Equivariant projection]
  \label{def:equivariant-projection}
  \lean{LinearMap.equivariantProjection}
  \mathlibok
  % [ML] LinearMap.equivariantProjection G π = (1/|G|) • ∑ g, π.conjugate g.
  \uses{def:sum-of-conjugates}
  Set $\bar\pi := |G|^{-1}\,\Sigma\pi$. The Mathlib name is
  \texttt{LinearMap.equivariantProjection G π} and pointwise
  \texttt{equivariantProjection\_apply}: $\bar\pi(v) =
  \frac{1}{|G|}\sum_g \pi^g(v)$.
\end{definition}

\begin{lemma}[$\bar\pi$ is ${k[G]}$-linear]
  \label{lem:bar-pi-equivariant}
  \lean{LinearMap.equivariantProjection}
  \mathlibok
  % [ML] The Mathlib definition LinearMap.equivariantProjection bundles
  %   k[G]-linearity into its type: it returns a (W →ₗ[MonoidAlgebra k G] V).
  \uses{def:equivariant-projection}
  The map $\bar\pi : W \to V$ is $k[G]$-linear: it commutes with the action
  of every group element.
\end{lemma}

\begin{proof}
  \leanok
  \uses{def:equivariant-projection}
  Bundled into the Mathlib definition: \texttt{LinearMap.equivariantProjection}
  has codomain \texttt{(W →ₗ[MonoidAlgebra k G] V)}. The proof in Mathlib is by
  reindexing the sum $\sum_g \pi^g$: conjugating by $h$ permutes the summands
  ($g \mapsto h^{-1} g$) and the prefactor $|G|^{-1}$ is fixed.
\end{proof}

\begin{lemma}[$\bar\pi$ retracts $i$]
  \label{lem:bar-pi-retracts}
  \lean{LinearMap.equivariantProjection_condition}
  \mathlibok
  % [ML] Mathlib.RepresentationTheory.Maschke:
  %   LinearMap.equivariantProjection_condition.
  \uses{def:equivariant-projection, lem:conjugate-of-i}
  If $i : V \to W$ is $k[G]$-linear and $\pi \circ i = \mathrm{id}_V$, then
  $\bar\pi \circ i = \mathrm{id}_V$.
\end{lemma}

\begin{proof}
  \leanok
  \uses{lem:conjugate-of-i}
  Compute, for $v \in V$:
  $\bar\pi(i(v)) = |G|^{-1}\sum_g \pi^g(i(v)) = |G|^{-1}\sum_g v
  = |G|^{-1}\cdot|G|\cdot v = v$,
  using Lemma~\ref{lem:conjugate-of-i} for the second equality and the
  invertibility of $|G|$ for the last. This is
  \texttt{LinearMap.equivariantProjection\_condition}.
\end{proof}

\section{Maschke's Theorem}

\begin{theorem}[Maschke: every ${k[G]}$-injection splits]
  \label{thm:maschke-split}
  \lean{MonoidAlgebra.exists_leftInverse_of_injective}
  \mathlibok
  % [ML] Mathlib.RepresentationTheory.Maschke:
  %   MonoidAlgebra.exists_leftInverse_of_injective.
  %   Hypotheses: [Field k] [Group G] [Finite G] [NeZero (Nat.card G : k)].
  \uses{lem:bar-pi-equivariant, lem:bar-pi-retracts}
  Let $V, W$ be $k[G]$-modules and $f : V \to W$ a $k[G]$-linear injection.
  Then $f$ admits a $k[G]$-linear left inverse: there exists
  $g : W \to V$ in $k[G]\text{-Mod}$ with $g \circ f = \mathrm{id}_V$.
\end{theorem}

\begin{proof}
  \leanok
  \uses{lem:bar-pi-equivariant, lem:bar-pi-retracts}
  Choose any $k$-linear retraction $\pi : W \to V$ of $f$ (it exists because
  $V \to W$ is an injection of $k$-vector spaces). Take
  $g := \bar\pi$ from Definition~\ref{def:equivariant-projection}; it is
  $k[G]$-linear (Lemma~\ref{lem:bar-pi-equivariant}) and retracts $f$
  (Lemma~\ref{lem:bar-pi-retracts}).  This is the proof of
  \texttt{MonoidAlgebra.exists\_leftInverse\_of\_injective}.
\end{proof}

\begin{theorem}[Maschke: every submodule has a complement]
  \label{thm:maschke-compl}
  \lean{MonoidAlgebra.Submodule.exists_isCompl}
  \mathlibok
  % [ML] Mathlib.RepresentationTheory.Maschke:
  %   MonoidAlgebra.Submodule.exists_isCompl.
  \uses{thm:maschke-split}
  For any $k[G]$-module $V$ and any $k[G]$-submodule $p \subseteq V$ there is
  a $k[G]$-submodule $q \subseteq V$ with $V = p \oplus q$ (i.e.\
  \texttt{IsCompl p q}).
\end{theorem}

\begin{proof}
  \leanok
  \uses{thm:maschke-split}
  Apply Theorem~\ref{thm:maschke-split} to the inclusion
  $p \hookrightarrow V$ to obtain a $k[G]$-linear retraction
  $g : V \to p$. Then $q := \ker g$ is a $k[G]$-submodule with $p \oplus q = V$.
\end{proof}

\begin{theorem}[Maschke: $V$ is semisimple as a ${k[G]}$-module]
  \label{thm:maschke-semisimple}
  \lean{MonoidAlgebra.Submodule.instIsSemisimpleModule}
  \mathlibok
  % [ML] Mathlib.RepresentationTheory.Maschke registers this as an instance
  %   IsSemisimpleModule (MonoidAlgebra k G) V.
  \uses{thm:maschke-compl}
  For every $k[G]$-module $V$ we have
  \texttt{IsSemisimpleModule (MonoidAlgebra k G) V}: every submodule has a
  complement, equivalently, $V$ is a direct sum of simple submodules.
\end{theorem}

\begin{proof}
  \leanok
  \uses{thm:maschke-compl}
  This is the Mathlib instance \texttt{MonoidAlgebra.Submodule.instIsSemisimpleModule},
  which packages Theorem~\ref{thm:maschke-compl} into the typeclass
  \texttt{IsSemisimpleModule}.
\end{proof}

\begin{corollary}[Every f.d.\ rep is a direct sum of irreducibles]
  \label{cor:fd-rep-decomposition}
  \lean{IsSemisimpleModule.exists_directSum_simple}
  % [~ML] Generic statement available for any IsSemisimpleModule on a
  %       Noetherian module; for finite-dimensional vector spaces this
  %       reduces to a finite direct sum.
  \uses{thm:maschke-semisimple, def:simple}
  Every $V : \texttt{FDRep k G}$ is isomorphic to a finite direct sum of simple
  objects of \texttt{FDRep k G}.
\end{corollary}

\begin{proof}
  \uses{thm:maschke-semisimple}
  By Theorem~\ref{thm:maschke-semisimple}, \texttt{V.ρ.asModule} is semisimple.
  A finitely generated semisimple module is a finite direct sum of simple
  modules; transporting along the equivalence
  \texttt{Rep.equivalenceModuleMonoidAlgebra} gives the statement in
  \texttt{FDRep k G}.
\end{proof}


\chapter{Schur's Lemma and Matrix-Coefficient Orthogonality}
\label{chap:schur}

\section{Schur's Lemma}

\begin{theorem}[Schur's Lemma — dimension form]
  \label{thm:schur-dim}
  \lean{FDRep.finrank_hom_simple_simple}
  \mathlibok
  % [ML] Mathlib.RepresentationTheory.FDRep:
  %   FDRep.finrank_hom_simple_simple :
  %     finrank k (V ⟶ W) = if Nonempty (V ≅ W) then 1 else 0,
  %   for V W : FDRep k G with [Simple V] [Simple W] [IsAlgClosed k].
  \uses{def:simple, def:intertwiner}
  Let $V, W : \texttt{FDRep k G}$ both be simple. Then
  \[
    \dim_k(V \to W)_G \;=\;
    \begin{cases} 1 & \text{if } V \cong W,\\ 0 & \text{otherwise.} \end{cases}
  \]
\end{theorem}

\begin{proof}
  \leanok
  \uses{def:simple, def:intertwiner}
  This is \texttt{FDRep.finrank\_hom\_simple\_simple}.  The proof packaged in
  Mathlib uses the more general \texttt{CategoryTheory.Schur} for an abelian
  $k$-linear category with simple objects and $k$ algebraically closed:
  any nonzero morphism $T : V \to W$ between simples has zero kernel and
  zero cokernel by simplicity of $V, W$, hence is an iso; for endomorphisms
  $T : V \to V$, since $k$ is algebraically closed $T$ has an eigenvalue
  $\lambda$, and $T - \lambda \mathrm{id}$ is a non-iso endomorphism of a simple
  object, hence zero.
\end{proof}

\begin{corollary}[Schur scalar form]
  \label{cor:schur-scalar}
  \lean{CategoryTheory.endRingEquiv}
  % [~ML] CategoryTheory.endRingEquiv (or End_simple_isAlgEquiv) gives
  %       End V ≃+* k for V : FDRep k G simple, under [IsAlgClosed k].
  %       Direct corollary of Theorem~\ref{thm:schur-dim} with V = W.
  \uses{thm:schur-dim}
  For $V : \texttt{FDRep k G}$ simple, every $G$-equivariant endomorphism of
  $V$ is a scalar multiple of $\mathrm{id}_V$. The ring $\mathrm{End}_G(V)$ is
  isomorphic to $k$.
\end{corollary}

\begin{proof}
  \uses{thm:schur-dim}
  By Theorem~\ref{thm:schur-dim} with $V = W$,
  $\dim_k \mathrm{End}_G V = 1$. Since $\mathrm{id}_V \in \mathrm{End}_G V$ is
  nonzero, it spans, so every endomorphism is a scalar.
\end{proof}

\section{The Symmetrisation \texorpdfstring{$A \mapsto |G|^{-1}\sum_g \rho(g) A \rho(g)^{-1}$}{A -> |G|\^{}-1 sum rho(g) A rho(g)\^{}-1}}

The following few lemmas isolate the averaging argument used in the
matrix-coefficient orthogonality calculation. They are entirely elementary
and do not yet appear as standalone Mathlib declarations.

\begin{definition}[$G$-symmetrisation of a linear map]
  \label{def:symmetrisation}
  \lean{Representation.symmetrise}
  % [NEW] Not in Mathlib in this form.  The construction is
  %   ρ.symmetrise σ A := |G|⁻¹ • ∑ g, σ g ∘ A ∘ ρ g⁻¹.
  \uses{def:representation, def:standing-hypotheses}
  Given $\rho : \mathrm{Representation}\ k\ G\ V$,
  $\sigma : \mathrm{Representation}\ k\ G\ W$ and a $k$-linear $A : V \to W$,
  define
  \[
    \mathrm{Sym}_{\rho,\sigma}(A) \;:=\; \frac{1}{|G|}\sum_{g \in G}
        \sigma(g)\, A\, \rho(g)^{-1} \;\in\; (V \to_{\!\ell[k]} W).
  \]
\end{definition}

\begin{lemma}[Symmetrisation lands in equivariant maps]
  \label{lem:symmetrise-equivariant}
  \lean{Representation.symmetrise_equivariant}
  % [NEW]
  \uses{def:symmetrisation, def:intertwiner}
  $\mathrm{Sym}_{\rho,\sigma}(A)$ is $G$-equivariant: it intertwines $\rho$ and
  $\sigma$.
\end{lemma}

\begin{proof}
  \uses{def:symmetrisation}
  Compute, for $h \in G$:
  \begin{align*}
    \sigma(h)\,\mathrm{Sym}(A) &= \tfrac{1}{|G|}\textstyle\sum_g \sigma(h)\sigma(g) A \rho(g)^{-1}
       = \tfrac{1}{|G|}\textstyle\sum_g \sigma(hg) A \rho(g)^{-1} \\
    &= \tfrac{1}{|G|}\textstyle\sum_{g'} \sigma(g') A \rho(h^{-1}g')^{-1}
       \quad (g' := hg) \\
    &= \tfrac{1}{|G|}\textstyle\sum_{g'} \sigma(g') A \rho(g')^{-1}\rho(h)
       = \mathrm{Sym}(A)\,\rho(h).
  \end{align*}
\end{proof}

\begin{lemma}[Symmetrisation between non-isomorphic simples is zero]
  \label{lem:symmetrise-nonzero-implies-iso}
  \lean{Representation.symmetrise_eq_zero_of_not_iso}
  % [NEW]
  \uses{lem:symmetrise-equivariant, thm:schur-dim, def:simple}
  If $V, W : \texttt{FDRep k G}$ are simple and not isomorphic, then
  $\mathrm{Sym}_{V.\rho, W.\rho}(A) = 0$ for every $k$-linear $A : V \to W$.
\end{lemma}

\begin{proof}
  \uses{lem:symmetrise-equivariant, thm:schur-dim}
  By Lemma~\ref{lem:symmetrise-equivariant}, $\mathrm{Sym}(A)$ is an
  intertwiner $V \to W$. By Theorem~\ref{thm:schur-dim}, the space of such
  intertwiners has dimension $0$ since $V \not\cong W$, hence
  $\mathrm{Sym}(A) = 0$.
\end{proof}

\begin{lemma}[Symmetrisation on a single simple is a scalar]
  \label{lem:symmetrise-self-scalar}
  \lean{Representation.symmetrise_self_eq_smul_id}
  % [NEW]
  \uses{lem:symmetrise-equivariant, cor:schur-scalar, def:simple}
  Let $V : \texttt{FDRep k G}$ be simple. For every $A : V \to V$ ($k$-linear),
  there is a unique scalar $c(A) \in k$ such that
  $\mathrm{Sym}_{V.\rho,V.\rho}(A) = c(A)\,\mathrm{id}_V$.  Furthermore
  \[
    c(A) \;=\; \frac{\mathrm{tr}(A)}{\dim_k V}.
  \]
\end{lemma}

\begin{proof}
  \uses{cor:schur-scalar, lem:symmetrise-equivariant}
  By Lemma~\ref{lem:symmetrise-equivariant} the symmetrisation is in
  $\mathrm{End}_G V$, which by Corollary~\ref{cor:schur-scalar} is one-dimensional
  spanned by $\mathrm{id}_V$. Hence $\mathrm{Sym}(A) = c(A)\mathrm{id}_V$ for some
  $c(A) \in k$. Take the $k$-trace of both sides: the left side has trace
  $\frac{1}{|G|}\sum_g \mathrm{tr}(\rho(g) A \rho(g)^{-1}) = \mathrm{tr}(A)$
  (cyclicity of trace), and the right side has trace $c(A)\,\dim_k V$.
  So $c(A) = \mathrm{tr}(A)/\dim_k V$.
\end{proof}

\section{Matrix-Coefficient Orthogonality (Schur II)}

\begin{theorem}[Matrix-coefficient orthogonality, single irrep]
  \label{thm:matrix-coeff-ortho-self}
  \lean{Representation.matrix_coeff_orthogonality_self}
  % [NEW] Standard consequence of Lemma~\ref{lem:symmetrise-self-scalar};
  %       not packaged in Mathlib.
  \uses{lem:symmetrise-self-scalar, def:simple}
  Let $V : \texttt{FDRep k G}$ be simple of dimension $d := \dim_k V$, and let
  $\{e_1, \ldots, e_d\}$ be a $k$-basis of $V$. Write $\rho_{ij}(g)$ for the
  $(i,j)$-th matrix coefficient of $V.\rho(g)$ in this basis. Then
  \[
    \frac{1}{|G|}\sum_{g \in G} \rho_{ij}(g)\, \rho_{kl}(g^{-1})
    \;=\; \frac{\delta_{il}\,\delta_{jk}}{d}.
  \]
\end{theorem}

\begin{proof}
  \uses{lem:symmetrise-self-scalar}
  Apply Lemma~\ref{lem:symmetrise-self-scalar} to the rank-one $k$-linear map
  $A := E_{lj}$ (the matrix unit sending $e_j$ to $e_l$, all other basis
  vectors to $0$). Then $\mathrm{tr}(A) = \delta_{lj}$, so
  $c(A) = \delta_{lj}/d$, and reading off the $(i,k)$-th matrix entry of
  $\mathrm{Sym}(A) = c(A)\mathrm{id}_V$ on both sides gives
  $\frac{1}{|G|}\sum_g \rho_{ij}(g)\rho_{kl}(g^{-1}) = \delta_{lj}\delta_{ik}/d$.
\end{proof}

\begin{theorem}[Matrix-coefficient orthogonality, distinct irreps]
  \label{thm:matrix-coeff-ortho-distinct}
  \lean{Representation.matrix_coeff_orthogonality_distinct}
  % [NEW] Counterpart of \ref{thm:matrix-coeff-ortho-self} for V ≠ W.
  \uses{lem:symmetrise-nonzero-implies-iso, def:simple}
  Let $V, W : \texttt{FDRep k G}$ be simple with $V \not\cong W$, and choose
  bases of $V$ and $W$. Then for all matrix indices,
  \[
    \frac{1}{|G|}\sum_{g \in G} \rho_{ij}(g)\, \sigma_{kl}(g^{-1}) \;=\; 0,
  \]
  where $\rho = V.\rho$ and $\sigma = W.\rho$.
\end{theorem}

\begin{proof}
  \uses{lem:symmetrise-nonzero-implies-iso}
  Apply Lemma~\ref{lem:symmetrise-nonzero-implies-iso} to any rank-one
  $A : V \to W$ and read off the result in coordinates.
\end{proof}


\chapter{Characters and Their Orthogonality}
\label{chap:characters}

\section{Characters}

\begin{definition}[Character of a representation]
  \label{def:character}
  \lean{FDRep.character}
  \mathlibok
  % [ML] Mathlib.RepresentationTheory.Character:
  %   FDRep.character V g = LinearMap.trace k V (V.ρ g).
  \uses{def:rep-categories}
  The \emph{character} of $V : \texttt{FDRep k G}$ is the function
  $\chi_V : G \to k$, $g \mapsto \mathrm{tr}_V(V.\rho(g))$.
\end{definition}

\begin{lemma}[$\chi_V(1) = \dim V$]
  \label{lem:char-one}
  \lean{FDRep.char_one}
  \mathlibok
  % [ML] FDRep.char_one : V.character 1 = finrank k V.
  \uses{def:character}
  $\chi_V(1) = \dim_k V$.
\end{lemma}

\begin{proof}
  \leanok
  \uses{def:character}
  $\chi_V(1) = \mathrm{tr}(V.\rho(1)) = \mathrm{tr}(\mathrm{id}_V) = \dim_k V$.
  This is \texttt{FDRep.char\_one}.
\end{proof}

\begin{lemma}[Characters are class functions]
  \label{lem:char-conj}
  \lean{FDRep.char_conj}
  \mathlibok
  % [ML] FDRep.char_conj : V.character (h * g * h⁻¹) = V.character g.
  \uses{def:character}
  $\chi_V(hgh^{-1}) = \chi_V(g)$ for all $g, h \in G$.
\end{lemma}

\begin{proof}
  \leanok
  \uses{def:character}
  Cyclicity of trace: $\mathrm{tr}(\rho(h)\rho(g)\rho(h)^{-1}) = \mathrm{tr}(\rho(g))$.
  Mathlib: \texttt{FDRep.char\_conj}.
\end{proof}

\begin{lemma}[Characters are isomorphism invariants]
  \label{lem:char-iso}
  \lean{FDRep.char_iso}
  \mathlibok
  % [ML] FDRep.char_iso : (V ≅ W) → V.character = W.character.
  \uses{def:character, def:rep-iso}
  If $V \cong W$ in \texttt{FDRep k G} then $\chi_V = \chi_W$.
\end{lemma}

\begin{proof}
  \leanok
  \uses{def:character}
  An isomorphism conjugates one action into the other; trace is conjugation-invariant.
  Mathlib: \texttt{FDRep.char\_iso}.
\end{proof}

\section{The Averaging Formula and Character Orthonormality}

\begin{theorem}[Average of $\chi_V$ counts invariants]
  \label{thm:avg-char-eq-invariants}
  \lean{FDRep.average_char_eq_finrank_invariants}
  \mathlibok
  % [ML] FDRep.average_char_eq_finrank_invariants:
  %   ⅟|G| • ∑ g, V.character g = finrank k (Representation.invariants V.ρ).
  \uses{def:character, def:subrep}
  For $V : \texttt{FDRep k G}$,
  \[
    \frac{1}{|G|}\sum_{g \in G} \chi_V(g) \;=\; \dim_k V^G.
  \]
\end{theorem}

\begin{proof}
  \leanok
  \uses{def:character, def:subrep}
  This is \texttt{FDRep.average\_char\_eq\_finrank\_invariants}.
  The standard proof: the operator $P := \frac{1}{|G|}\sum_g \rho(g)$ is a
  $G$-equivariant idempotent ($P^2 = P$, since left-translation by $h$
  permutes the summands and the prefactor is invariant) whose image is exactly
  $V^G$. The trace of an idempotent endomorphism of a finite-dimensional
  vector space is the dimension of its image, giving
  $\frac{1}{|G|}\sum_g\chi_V(g) = \mathrm{tr}(P) = \dim V^G$.
\end{proof}

\begin{theorem}[Hom-space dimension via characters]
  \label{thm:hom-space-dim-via-char}
  \lean{Representation.card_inv_mul_sum_char_mul_char_eq_finrank}
  \mathlibok
  % [ML] Representation.card_inv_mul_sum_char_mul_char_eq_finrank:
  %   ⅟|G| • ∑ g, V.character g * W.character g⁻¹ = finrank k (V ⟶ W).
  \uses{thm:avg-char-eq-invariants, def:character, def:intertwiner}
  For $V, W : \texttt{FDRep k G}$,
  \[
    \frac{1}{|G|}\sum_{g \in G} \chi_V(g)\, \chi_W(g^{-1})
    \;=\; \dim_k(V \to W)_G.
  \]
\end{theorem}

\begin{proof}
  \leanok
  \uses{thm:avg-char-eq-invariants}
  Apply Theorem~\ref{thm:avg-char-eq-invariants} to the representation
  $\mathrm{Hom}_k(V, W)$ with $G$-action $g \cdot T := \rho_W(g)\, T\, \rho_V(g)^{-1}$.
  Its character at $g$ equals $\chi_V(g^{-1})\chi_W(g)$, and its
  $G$-invariants are exactly $(V \to W)_G$. Mathlib:
  \texttt{Representation.card\_inv\_mul\_sum\_char\_mul\_char\_eq\_finrank}.
\end{proof}

\begin{theorem}[First Schur orthogonality of characters]
  \label{thm:char-ortho}
  \lean{FDRep.char_orthonormal}
  \mathlibok
  % [ML] FDRep.char_orthonormal:
  %   ⅟|G| • ∑ g, V.character g * W.character g⁻¹
  %     = if Nonempty (V ≅ W) then 1 else 0,
  %   for V W : FDRep k G simple, [IsAlgClosed k].
  \uses{thm:hom-space-dim-via-char, thm:schur-dim}
  Let $V, W : \texttt{FDRep k G}$ both be simple. Then
  \[
    \langle \chi_V, \chi_W \rangle_G \;=\;
    \begin{cases} 1 & \text{if } V \cong W,\\ 0 & \text{otherwise.}\end{cases}
  \]
\end{theorem}

\begin{proof}
  \leanok
  \uses{thm:hom-space-dim-via-char, thm:schur-dim}
  Combine Theorems \ref{thm:hom-space-dim-via-char} and \ref{thm:schur-dim}:
  the left-hand side equals $\dim_k(V \to W)_G$ which equals $1$ or $0$
  depending on whether $V \cong W$. Mathlib: \texttt{FDRep.char\_orthonormal}.
\end{proof}


\chapter{Isotypic Decomposition}
\label{chap:isotypic}

\section{The Indexing Set $\hat G$}

\begin{definition}[Set of irreducible classes]
  \label{def:Ghat}
  \lean{FDRep.IrreducibleClass}
  % [NEW] Mathlib does not currently package the set of isomorphism classes
  %       of finite-dimensional simple representations as a standalone type.
  %       For finite groups this set is finite (Theorem~\ref{thm:Ghat-finite}).
  \uses{def:simple, def:rep-iso}
  Let $\hat G$ denote the set of isomorphism classes of simple objects of
  \texttt{FDRep k G}. For $\xi \in \hat G$ we choose a representative
  $V_\xi : \texttt{FDRep k G}$ with $\xi = [V_\xi]$ and write
  $d_\xi := \dim_k V_\xi$.
\end{definition}

\begin{theorem}[Finiteness of $\hat G$]
  \label{thm:Ghat-finite}
  \lean{FDRep.IrreducibleClass.fintype}
  % [NEW] Follows from the regular representation being finite-dimensional
  %       and decomposing into all irreps; concretely a consequence of the
  %       Peter–Weyl theorem (Theorem~\ref{thm:peter-weyl}). Can also be
  %       proved directly from $|\hat G| \le \dim k[G] = |G|$.
  \uses{def:Ghat}
  $\hat G$ is finite, with $|\hat G| \le |G|$. (We will see in
  Theorem~\ref{thm:char-basis} that in fact $|\hat G|$ equals the number of
  conjugacy classes of $G$.)
\end{theorem}

\begin{proof}
  \uses{thm:peter-weyl, def:Ghat}
  Every simple object embeds into the regular representation $k[G]$ (a
  consequence of Peter–Weyl: see Theorem~\ref{thm:peter-weyl} and the remark
  there that every irrep appears in $k[G]$ with multiplicity $d_\xi \ge 1$),
  and $k[G]$ is $|G|$-dimensional, so $|\hat G| \le |G|$.
\end{proof}

\section{Isotypic Components}

\begin{definition}[Isotypic component]
  \label{def:isotypic-component}
  \lean{Representation.isotypicComponent}
  % [NEW] Not in Mathlib.
  \uses{def:Ghat, def:subrep}
  For $\rho$ a representation on $V$ and $\xi \in \hat G$, the
  \emph{$\xi$-isotypic component} of $V$ is
  \[
    V_{(\xi)} \;:=\; \sum \big\{\, W \subseteq V : W\ \text{is a $G$-subspace
    isomorphic to } V_\xi\,\big\},
  \]
  the sum being taken inside $V$. It is itself a $G$-subrepresentation.
\end{definition}

\begin{lemma}[Isotypic components are pairwise disjoint as subspaces]
  \label{lem:isotypic-disjoint}
  \lean{Representation.isotypicComponent_disjoint}
  % [NEW]
  \uses{def:isotypic-component, thm:schur-dim}
  For $\xi \ne \xi'$ in $\hat G$, $V_{(\xi)} \cap V_{(\xi')} = 0$ inside $V$.
\end{lemma}

\begin{proof}
  \uses{thm:schur-dim, def:isotypic-component, thm:maschke-semisimple}
  Write $W := V_{(\xi)} \cap V_{(\xi')}$, a $G$-subrepresentation of $V$.
  By Maschke (Theorem~\ref{thm:maschke-semisimple}) $W$ is semisimple,
  i.e.\ a direct sum of simples. Each simple summand sits in both
  $V_{(\xi)}$ and $V_{(\xi')}$, so by the definition of isotypic component
  is isomorphic to both $V_\xi$ and $V_{\xi'}$, contradicting $\xi \ne \xi'$.
  Hence there are no simple summands and $W = 0$.
\end{proof}

\begin{lemma}[Multiplicity space]
  \label{lem:multiplicity-space}
  \lean{Representation.multiplicitySpace}
  % [NEW]
  \uses{def:isotypic-component, thm:schur-dim, def:simple}
  For $V : \texttt{FDRep k G}$ and $\xi \in \hat G$, define
  $M_\xi := (V_\xi \to V)_G$ (the $k$-vector space of $G$-equivariant maps
  $V_\xi \to V$). Then $\dim_k M_\xi$ is the \emph{multiplicity} $m_\xi$ of
  $V_\xi$ in $V$, and there is a canonical $G$-equivariant isomorphism
  \[
    \mathrm{ev}_\xi : V_\xi \otimes_k M_\xi \xrightarrow{\ \sim\ } V_{(\xi)},
    \qquad v \otimes \phi \longmapsto \phi(v).
  \]
  In particular $V_{(\xi)} \cong V_\xi^{\oplus m_\xi}$ and
  $\dim_k V_{(\xi)} = m_\xi \cdot d_\xi$.
\end{lemma}

\begin{proof}
  \uses{thm:schur-dim, thm:maschke-semisimple, def:isotypic-component}
  Equivariance of $\mathrm{ev}_\xi$ is direct.  Surjectivity: by
  Theorem~\ref{thm:maschke-semisimple}, $V_{(\xi)}$ is a finite direct sum
  $\bigoplus_{\alpha} W_\alpha$ of simples isomorphic to $V_\xi$; pick
  isomorphisms $\phi_\alpha : V_\xi \to W_\alpha \subseteq V$, then any
  $w \in V_{(\xi)}$ writes as $\sum_\alpha \phi_\alpha(v_\alpha) =
  \mathrm{ev}_\xi\big(\sum_\alpha v_\alpha \otimes \phi_\alpha\big)$.
  Injectivity: an element of $V_\xi \otimes M_\xi$ in the kernel decomposes
  as $\sum v_\alpha \otimes \phi_\alpha$ for $\phi_\alpha$ a basis of
  $M_\xi$; the requirement $\sum \phi_\alpha(v_\alpha) = 0$ together with
  Theorem~\ref{thm:schur-dim} (which says the $\phi_\alpha$ have linearly
  independent images among simples) forces every $v_\alpha = 0$.
  The dimension count follows.
\end{proof}

\begin{theorem}[Isotypic decomposition]
  \label{thm:isotypic-decomposition}
  \lean{Representation.isotypicDecomposition}
  % [NEW]
  \uses{lem:isotypic-disjoint, lem:multiplicity-space, thm:maschke-semisimple}
  For every $V : \texttt{FDRep k G}$, the inclusions $V_{(\xi)} \hookrightarrow V$
  assemble into a $G$-equivariant isomorphism
  \[
    \bigoplus_{\xi \in \hat G} V_{(\xi)} \xrightarrow{\ \sim\ } V,
  \]
  and consequently
  \[
    V \;\cong\; \bigoplus_{\xi \in \hat G} V_\xi^{\oplus m_\xi},
    \qquad \dim_k V \;=\; \sum_{\xi \in \hat G} m_\xi\, d_\xi.
  \]
\end{theorem}

\begin{proof}
  \uses{lem:isotypic-disjoint, lem:multiplicity-space, thm:maschke-semisimple}
  By Theorem~\ref{thm:maschke-semisimple} write $V = \bigoplus_\alpha W_\alpha$
  with each $W_\alpha$ simple. Group the summands by their isomorphism
  class $\xi \in \hat G$; the sum of all $W_\alpha$ in class $\xi$ is exactly
  $V_{(\xi)}$ (it is contained in $V_{(\xi)}$ and contains every simple summand
  of $V$ isomorphic to $V_\xi$). The $V_{(\xi)}$ jointly span $V$ by
  construction, and pairwise intersect trivially by
  Lemma~\ref{lem:isotypic-disjoint}; hence they form an internal direct sum
  decomposition of $V$. Combine with the description from
  Lemma~\ref{lem:multiplicity-space}.
\end{proof}


\chapter{The Peter–Weyl Theorem for Finite Groups}
\label{chap:peter-weyl}

\section{The Regular Representation}

\begin{definition}[Left regular representation]
  \label{def:regular}
  \lean{Representation.leftRegular, Rep.ofMulAction}
  \mathlibok
  % [ML] Mathlib.RepresentationTheory.Basic: Representation.leftRegular k G;
  %      Mathlib.RepresentationTheory.Rep: Rep.ofMulAction k G G — the
  %      bundled version on k[G] = (G →₀ k).
  \uses{def:representation, def:rep-categories}
  The \emph{left regular representation} is
  $\lambda : G \to (k[G] \to_{\!\ell[k]} k[G])$ given by
  $\lambda(g)(\delta_h) = \delta_{gh}$ on the basis $\{\delta_h : h \in G\}$.
  Mathlib: \texttt{Representation.leftRegular k G} (and its bundled form
  \texttt{Rep.ofMulAction k G G : Rep k G}).
\end{definition}

\begin{lemma}[Character of the regular representation]
  \label{lem:char-regular}
  \lean{Representation.character_leftRegular}
  % [~ML] The pointwise formula is folklore; not currently a single named
  %       Mathlib lemma, but trace of the permutation matrix shows
  %       $\chi_{k[G]}(g) = |\{h : g \cdot h = h\}| = |G|\cdot[g = 1]$.
  \uses{def:regular, def:character}
  $\chi_{k[G]}(g) = |G|$ if $g = 1$, and $0$ otherwise.
\end{lemma}

\begin{proof}
  \uses{def:regular, def:character}
  The matrix of $\lambda(g)$ in the basis $\{\delta_h\}_{h \in G}$ is the
  permutation matrix of left translation by $g$. Its trace is the number of
  fixed points of $h \mapsto gh$, which is $|G|$ if $g = 1$ and $0$ otherwise
  (a free $G$-action on $G$ has no fixed points unless $g = 1$).
\end{proof}

\section{Multiplicities in $k[G]$}

\begin{lemma}[Multiplicity of $V_\xi$ in ${k[G]}$]
  \label{lem:mult-in-regular}
  \lean{Representation.multiplicity_in_leftRegular}
  % [NEW]
  \uses{lem:char-regular, thm:hom-space-dim-via-char, def:Ghat}
  Let $\xi \in \hat G$ and $m_\xi$ denote the multiplicity of $V_\xi$ in
  $k[G]$ (Lemma~\ref{lem:multiplicity-space}). Then $m_\xi = d_\xi$.
\end{lemma}

\begin{proof}
  \uses{lem:char-regular, thm:hom-space-dim-via-char, lem:multiplicity-space, lem:char-one}
  By Lemma~\ref{lem:multiplicity-space}, $m_\xi = \dim_k(V_\xi \to k[G])_G$.
  By Theorem~\ref{thm:hom-space-dim-via-char},
  \[
    \dim_k(V_\xi \to k[G])_G \;=\; \frac{1}{|G|}\sum_g \chi_{V_\xi}(g)\,\chi_{k[G]}(g^{-1}).
  \]
  Using Lemma~\ref{lem:char-regular} only the $g = 1$ term survives, giving
  $\frac{1}{|G|} \cdot \chi_{V_\xi}(1)\cdot|G| = \chi_{V_\xi}(1) = d_\xi$
  by Lemma~\ref{lem:char-one}.
\end{proof}

\section{The Peter–Weyl Decomposition}

\begin{theorem}[Peter–Weyl: decomposition of the regular representation]
  \label{thm:peter-weyl}
  \lean{PeterWeyl.regular_decomposition}
  % [NEW] Specialisation of Theorem~\ref{thm:isotypic-decomposition}
  %       to V = k[G], using Lemma~\ref{lem:mult-in-regular}.
  \uses{thm:isotypic-decomposition, lem:mult-in-regular}
  As $k$-linear $G$-representations,
  \[
    k[G] \;\cong\; \bigoplus_{\xi \in \hat G} V_\xi^{\oplus d_\xi}.
  \]
  In particular every irreducible representation appears in $k[G]$ with
  multiplicity equal to its dimension.
\end{theorem}

\begin{proof}
  \uses{thm:isotypic-decomposition, lem:mult-in-regular}
  Apply Theorem~\ref{thm:isotypic-decomposition} to $V := k[G]$, using
  Lemma~\ref{lem:mult-in-regular} to identify each multiplicity $m_\xi$ with
  $d_\xi$.
\end{proof}

\begin{corollary}[Sum of squares of dimensions]
  \label{cor:sum-of-squares}
  \lean{PeterWeyl.sum_sq_dim_eq_card}
  % [NEW] Immediate from Theorem~\ref{thm:peter-weyl} by counting dimensions.
  \uses{thm:peter-weyl}
  $\displaystyle |G| \;=\; \sum_{\xi \in \hat G} d_\xi^{\,2}$.
\end{corollary}

\begin{proof}
  \uses{thm:peter-weyl}
  Take $\dim_k$ of both sides of Theorem~\ref{thm:peter-weyl}: the left
  side is $|G|$, the right side is $\sum_\xi d_\xi \cdot d_\xi = \sum_\xi d_\xi^2$.
\end{proof}

\section{Two-Sided Wedderburn Form (optional)}

\begin{theorem}[Wedderburn form of the group algebra]
  \label{thm:wedderburn}
  \lean{PeterWeyl.groupAlgebra_iso_endomorphism_product}
  % [NEW] Strengthens Theorem~\ref{thm:peter-weyl} to a ring isomorphism.
  %       Mathlib has IsSemisimpleRing for k[G] under our hypotheses, but
  %       does not yet package the explicit Wedderburn decomposition
  %       k[G] ≅ ∏ End_k(V_ξ) as a ring isomorphism.
  \uses{thm:peter-weyl, cor:schur-scalar}
  As $k$-algebras (equivalently, as two-sided $k[G]$-modules) there is an
  isomorphism
  \[
    k[G] \;\cong\; \prod_{\xi \in \hat G} \mathrm{End}_k(V_\xi).
  \]
\end{theorem}

\begin{proof}
  \uses{thm:peter-weyl, cor:schur-scalar}
  The $k$-algebra map $\Phi : k[G] \to \prod_\xi \mathrm{End}_k(V_\xi)$
  sending $\sum_g a_g g \mapsto (\sum_g a_g\,V_\xi.\rho(g))_\xi$ is well-defined
  and an algebra homomorphism by definition of $V_\xi.\rho$ as a monoid hom.
  Surjectivity: under the equivalence
  \texttt{Rep.equivalenceModuleMonoidAlgebra}, $\Phi$ corresponds to the
  natural map
  $k[G] \to \mathrm{End}_{k[G]}\bigl(\bigoplus_\xi V_\xi^{\oplus d_\xi}\bigr)
  \cong \prod_\xi \mathrm{End}_k(V_\xi)^{\mathrm{op}}$
  (the second isomorphism uses Corollary~\ref{cor:schur-scalar} and the
  Schur identification $\mathrm{End}_{k[G]}(V_\xi^{\oplus d_\xi}) \cong
  M_{d_\xi}(k)$). Injectivity: source and target have equal dimension over $k$,
  namely $|G| = \sum_\xi d_\xi^2$ (Corollary~\ref{cor:sum-of-squares}); since
  $\Phi$ is surjective and source and target are finite-dimensional with
  matching dimension, $\Phi$ is bijective.
\end{proof}


\chapter{Multiplicity Formula and Character Theory Consequences}
\label{chap:character-consequences}

\section{Multiplicity from Characters}

\begin{theorem}[Multiplicity formula]
  \label{thm:multiplicity-formula}
  \lean{Representation.multiplicity_eq_inner_char}
  % [NEW] Combines Lemma~\ref{lem:multiplicity-space} with Theorem~\ref{thm:hom-space-dim-via-char}.
  \uses{lem:multiplicity-space, thm:hom-space-dim-via-char, def:character}
  For any $V : \texttt{FDRep k G}$ and $\xi \in \hat G$, the multiplicity
  $m_\xi$ of $V_\xi$ in $V$ satisfies
  \[
    m_\xi \;=\; \langle \chi_V, \chi_{V_\xi}\rangle_G
    \;=\; \frac{1}{|G|}\sum_{g \in G}\chi_V(g)\,\chi_{V_\xi}(g^{-1}).
  \]
\end{theorem}

\begin{proof}
  \uses{lem:multiplicity-space, thm:hom-space-dim-via-char}
  By Lemma~\ref{lem:multiplicity-space}, $m_\xi = \dim_k(V_\xi \to V)_G$.
  By Theorem~\ref{thm:hom-space-dim-via-char},
  $\dim_k(V_\xi \to V)_G = |G|^{-1}\sum_g \chi_{V_\xi}(g)\chi_V(g^{-1})$.
  Replacing $g \mapsto g^{-1}$ (a bijection of $G$) shows this equals
  $|G|^{-1}\sum_g \chi_V(g)\chi_{V_\xi}(g^{-1}) = \langle\chi_V,\chi_{V_\xi}\rangle_G$.
\end{proof}

\begin{corollary}[Characters separate isomorphism classes]
  \label{cor:char-separates}
  \lean{FDRep.iso_iff_character_eq}
  % [NEW] Immediate from Theorem~\ref{thm:multiplicity-formula}.
  \uses{thm:multiplicity-formula, lem:char-iso, thm:isotypic-decomposition}
  For $V, W : \texttt{FDRep k G}$, $V \cong W$ iff $\chi_V = \chi_W$.
\end{corollary}

\begin{proof}
  \uses{thm:multiplicity-formula, lem:char-iso, thm:isotypic-decomposition}
  ($\Rightarrow$) is Lemma~\ref{lem:char-iso}. ($\Leftarrow$):
  by Theorem~\ref{thm:multiplicity-formula}, $\chi_V = \chi_W$ implies the
  multiplicities $m_\xi$ for $V$ and $W$ coincide; by
  Theorem~\ref{thm:isotypic-decomposition} the multiplicities determine the
  isomorphism class, so $V \cong W$.
\end{proof}

\begin{corollary}[Irreducibility test]
  \label{cor:irred-test}
  \lean{FDRep.irreducible_iff_inner_self_eq_one}
  % [NEW]
  \uses{thm:multiplicity-formula, def:simple}
  $V : \texttt{FDRep k G}$ is simple iff $\langle \chi_V, \chi_V\rangle_G = 1$.
\end{corollary}

\begin{proof}
  \uses{thm:multiplicity-formula, thm:isotypic-decomposition}
  Write $V \cong \bigoplus_\xi V_\xi^{\oplus m_\xi}$. Then
  $\chi_V = \sum_\xi m_\xi \chi_{V_\xi}$ and by orthonormality
  (Theorem~\ref{thm:char-ortho}),
  $\langle \chi_V, \chi_V\rangle_G = \sum_\xi m_\xi^2$, an integer.
  This equals $1$ iff exactly one $m_\xi$ is $1$ and all others are $0$, iff
  $V \cong V_\xi$ is simple.
\end{proof}

\section{Class Functions and Their Span}

\begin{definition}[Space of class functions]
  \label{def:class-functions}
  \lean{ClassFunction}
  % [~ML] Easy local definition; not packaged as a standalone object in
  %       Mathlib but the underlying notion (conjugation invariance) is
  %       standard.
  \uses{def:standing-hypotheses}
  A \emph{class function} on $G$ is a function $f : G \to k$ with
  $f(hgh^{-1}) = f(g)$ for all $g, h$. Class functions form a $k$-subspace
  $\mathrm{Cl}(G) \subseteq (G \to k)$ whose dimension equals the number of
  conjugacy classes of $G$ (with delta functions on conjugacy classes as a
  basis).
\end{definition}

\begin{lemma}[$\dim_k \mathrm{Cl}(G) = $ number of conjugacy classes]
  \label{lem:dim-cl-G}
  \lean{ClassFunction.finrank_eq_conjClasses}
  % [~ML] Mathlib has ConjClasses; the dimension count is straightforward.
  \uses{def:class-functions}
  $\dim_k \mathrm{Cl}(G) = |\mathrm{ConjClasses}\ G|$.
\end{lemma}

\begin{proof}
  \uses{def:class-functions}
  A class function is determined by, and freely chosen on, each conjugacy
  class. The indicator functions of conjugacy classes form a basis.
\end{proof}

\begin{lemma}[Characters lie in $\mathrm{Cl}(G)$]
  \label{lem:char-in-cl}
  \lean{FDRep.character_mem_classFunction}
  % [~ML] Direct from Lemma~\ref{lem:char-conj}.
  \uses{def:character, def:class-functions, lem:char-conj}
  For every $V : \texttt{FDRep k G}$, $\chi_V \in \mathrm{Cl}(G)$.
\end{lemma}

\begin{proof}
  \leanok
  \uses{lem:char-conj}
  Lemma~\ref{lem:char-conj}.
\end{proof}

\section{Irreducible Characters Span $\mathrm{Cl}(G)$}

The next theorem is the second main result of finite-group character theory
after the Peter–Weyl decomposition. We isolate the spanning argument as a
separate lemma to enable independent formalisation.

\begin{lemma}[The projection $T_f$ associated to $f \in \mathrm{Cl}(G)$]
  \label{lem:Tf-projection}
  \lean{ClassFunction.toEndomorphism}
  % [NEW]
  \uses{def:class-functions, def:representation}
  For $f \in \mathrm{Cl}(G)$ and $\rho$ a representation of $G$ on $V$, define
  \[
    T_f^\rho \;:=\; \frac{1}{|G|}\sum_{g \in G} f(g)\,\rho(g) \;\in\;
    (V \to_{\!\ell[k]} V).
  \]
  Then $T_f^\rho$ is $G$-equivariant: $\rho(h)\, T_f^\rho = T_f^\rho\,\rho(h)$
  for every $h$.
\end{lemma}

\begin{proof}
  \uses{def:class-functions}
  Compute $\rho(h) T_f^\rho \rho(h)^{-1} = |G|^{-1}\sum_g f(g)\rho(hgh^{-1})
  = |G|^{-1}\sum_{g'}f(h^{-1}g'h)\rho(g')$ ($g' := hgh^{-1}$)
  $= |G|^{-1}\sum_{g'}f(g')\rho(g')$ (using $f$ is a class function)
  $= T_f^\rho$.
\end{proof}

\begin{lemma}[Schur scalar of $T_f$ on a simple]
  \label{lem:Tf-on-simple}
  \lean{ClassFunction.toEndomorphism_simple_smul_id}
  % [NEW]
  \uses{lem:Tf-projection, cor:schur-scalar, def:character, def:Ghat}
  For $\xi \in \hat G$, $T_f^{V_\xi} = \mu_\xi(f)\,\mathrm{id}_{V_\xi}$ where
  \[
    \mu_\xi(f) \;=\; \frac{1}{d_\xi \cdot |G|}\sum_{g \in G} f(g)\,\chi_{V_\xi}(g).
  \]
  In particular $T_f^{V_\xi} = 0$ iff $\sum_g f(g)\chi_{V_\xi}(g) = 0$.
\end{lemma}

\begin{proof}
  \uses{lem:Tf-projection, cor:schur-scalar}
  By Lemma~\ref{lem:Tf-projection} $T_f^{V_\xi}$ is in $\mathrm{End}_G(V_\xi)$,
  which by Corollary~\ref{cor:schur-scalar} is one-dimensional spanned by
  $\mathrm{id}_{V_\xi}$. So $T_f^{V_\xi} = \mu_\xi(f)\mathrm{id}_{V_\xi}$.
  Take $k$-trace of both sides:
  $\mathrm{tr}(T_f^{V_\xi}) = |G|^{-1}\sum_g f(g)\chi_{V_\xi}(g)$ on the left
  and $\mu_\xi(f)\cdot d_\xi$ on the right; solve for $\mu_\xi(f)$.
\end{proof}

\begin{lemma}[$f$ orthogonal to all $\chi_{V_\xi}$ implies $T_f^V = 0$ for every $V$]
  \label{lem:f-ortho-implies-Tf-zero}
  \lean{ClassFunction.toEndomorphism_eq_zero_of_orthogonal}
  % [NEW]
  \uses{lem:Tf-on-simple, thm:isotypic-decomposition}
  Let $f \in \mathrm{Cl}(G)$. If $\sum_g f(g)\chi_{V_\xi}(g) = 0$ for every
  $\xi \in \hat G$, then for every $V : \texttt{FDRep k G}$, $T_f^V = 0$.
\end{lemma}

\begin{proof}
  \uses{lem:Tf-on-simple, thm:isotypic-decomposition}
  By Theorem~\ref{thm:isotypic-decomposition}, $V$ is a direct sum of simples
  $V_\xi$, and $T_f^V$ acts diagonally on this decomposition by
  Lemma~\ref{lem:Tf-projection}. By Lemma~\ref{lem:Tf-on-simple} the
  hypothesis implies $T_f^{V_\xi} = 0$ on each summand, hence $T_f^V = 0$.
\end{proof}

\begin{lemma}[Applying $T_f$ on the regular representation to $\delta_1$]
  \label{lem:Tf-regular-delta-one}
  \lean{ClassFunction.toEndomorphism_regular_delta_one}
  % [NEW]
  \uses{lem:Tf-projection, def:regular}
  Let $\rho = \lambda$ be the regular representation on $V = k[G]$, and let
  $\delta_1 \in k[G]$ denote the basis vector at $1 \in G$. Then
  $T_f^\lambda(\delta_1) = \frac{1}{|G|}\sum_{g \in G} f(g)\,\delta_g
  = \frac{1}{|G|}\, f$, where on the right we identify $f \in (G \to k)$
  with the element $\sum_g f(g)\delta_g \in k[G]$.
\end{lemma}

\begin{proof}
  \leanok
  \uses{lem:Tf-projection, def:regular}
  Compute $T_f^\lambda(\delta_1) = |G|^{-1}\sum_g f(g)\,\lambda(g)\delta_1
  = |G|^{-1}\sum_g f(g)\,\delta_g$. Identifying $\delta_g$ with the basis
  function $\mathbf{1}_{\{g\}} \in (G \to k)$ and recombining gives
  $|G|^{-1}f$.
\end{proof}

\begin{theorem}[Irreducible characters span $\mathrm{Cl}(G)$]
  \label{thm:char-span}
  \lean{FDRep.span_irreducibleCharacters_eq_top}
  % [NEW]
  \uses{lem:f-ortho-implies-Tf-zero, lem:Tf-regular-delta-one, lem:char-in-cl}
  The irreducible characters $\{\chi_{V_\xi}\}_{\xi \in \hat G}$ span
  $\mathrm{Cl}(G)$.
\end{theorem}

\begin{proof}
  \uses{lem:f-ortho-implies-Tf-zero, lem:Tf-regular-delta-one}
  Suppose $f \in \mathrm{Cl}(G)$ satisfies $\sum_g f(g)\chi_{V_\xi}(g) = 0$
  for every $\xi \in \hat G$ (this is equivalent to $\langle f, \chi_{V_\xi^*}\rangle_G = 0$
  for every $\xi$, where $V_\xi^*$ is the dual representation; since
  $\xi \mapsto \xi^*$ permutes $\hat G$, this is the same as orthogonality
  to all irreducible characters in $\langle\cdot,\cdot\rangle_G$).
  By Lemma~\ref{lem:f-ortho-implies-Tf-zero} applied to the regular
  representation, $T_f^\lambda = 0$. Apply this to $\delta_1$:
  Lemma~\ref{lem:Tf-regular-delta-one} gives $0 = T_f^\lambda(\delta_1)
  = |G|^{-1}\,f$ in $k[G]$, hence $f = 0$ in $k^G$. So the orthogonal
  complement of the irreducible characters in $\mathrm{Cl}(G)$ is zero, i.e.\
  they span.
\end{proof}

\begin{theorem}[Irreducible characters are an orthonormal basis of $\mathrm{Cl}(G)$]
  \label{thm:char-basis}
  \lean{FDRep.irreducibleCharacters_orthonormalBasis}
  % [NEW] Packages Theorems \ref{thm:char-ortho} and \ref{thm:char-span}
  %       as a Basis (or OrthonormalBasis if k = ℂ).
  \uses{thm:char-ortho, thm:char-span, lem:dim-cl-G, def:Ghat}
  $\{\chi_{V_\xi}\}_{\xi \in \hat G}$ is an orthonormal basis of
  $\mathrm{Cl}(G)$ for the inner product $\langle\cdot,\cdot\rangle_G$.
  In particular,
  \[
    |\hat G| \;=\; \dim_k \mathrm{Cl}(G) \;=\; |\mathrm{ConjClasses}\ G|.
  \]
\end{theorem}

\begin{proof}
  \uses{thm:char-ortho, thm:char-span, lem:dim-cl-G}
  Orthonormality is Theorem~\ref{thm:char-ortho}, hence the
  $\chi_{V_\xi}$ are linearly independent. Spanning is
  Theorem~\ref{thm:char-span}. Hence they form a basis. The cardinality
  identity $|\hat G| = |\mathrm{ConjClasses}\ G|$ follows by comparing
  $\dim_k \mathrm{Cl}(G)$ via this basis with Lemma~\ref{lem:dim-cl-G}.
\end{proof}


\chapter{Summary of Mathlib Status}
\label{chap:status}

\paragraph{Already in Mathlib (\texttt{[ML]}).}
The following declarations underpin the entire blueprint:
\texttt{Representation}, \texttt{Representation.asModule}, \texttt{Rep},
\texttt{FDRep}, \texttt{CategoryTheory.Simple},
\texttt{Representation.invariants},
\texttt{LinearMap.conjugate}, \texttt{LinearMap.conjugate\_i},
\texttt{LinearMap.sumOfConjugatesEquivariant},
\texttt{LinearMap.equivariantProjection},
\texttt{LinearMap.equivariantProjection\_condition},
\texttt{MonoidAlgebra.exists\_leftInverse\_of\_injective},
\texttt{MonoidAlgebra.Submodule.exists\_isCompl},
\texttt{MonoidAlgebra.Submodule.instIsSemisimpleModule},
\texttt{FDRep.finrank\_hom\_simple\_simple} (Schur),
\texttt{FDRep.character}, \texttt{FDRep.char\_one},
\texttt{FDRep.char\_conj}, \texttt{FDRep.char\_iso},
\texttt{FDRep.char\_orthonormal},
\texttt{FDRep.average\_char\_eq\_finrank\_invariants},
\texttt{Representation.card\_inv\_mul\_sum\_char\_mul\_char\_eq\_finrank},
\texttt{Representation.leftRegular}, \texttt{Rep.ofMulAction}.

\paragraph{Close analogues (\texttt{[\textasciitilde ML]}).}
The space \texttt{ClassFunction} and the dimension count
$\dim \mathrm{Cl}(G) = |\mathrm{ConjClasses}\ G|$ are not packaged but
trivial. The character of the regular representation
($\chi_{k[G]} = |G|\cdot\mathbf{1}_{\{1\}}$) is folklore; trivial from the
permutation-matrix description of \texttt{Rep.ofMulAction}.

\paragraph{Newly required (\texttt{[NEW]}).}
The contributions to Mathlib that this blueprint amounts to are:
the inner product $\langle\cdot,\cdot\rangle_G$
(Definition~\ref{def:inner-prod}), the symmetrisation
(Definition~\ref{def:symmetrisation}) and matrix-coefficient orthogonality
(Theorems \ref{thm:matrix-coeff-ortho-self}, \ref{thm:matrix-coeff-ortho-distinct}),
the indexing set $\hat G$ and its finiteness, the isotypic component and
multiplicity space (Definition~\ref{def:isotypic-component},
Lemma~\ref{lem:multiplicity-space}), the isotypic decomposition
(Theorem~\ref{thm:isotypic-decomposition}),
the Peter–Weyl decomposition (Theorem~\ref{thm:peter-weyl}) and its
sum-of-squares corollary, the optional Wedderburn form
(Theorem~\ref{thm:wedderburn}), the multiplicity formula
(Theorem~\ref{thm:multiplicity-formula}) and its corollaries, and the spanning
of $\mathrm{Cl}(G)$ by irreducible characters
(Theorem~\ref{thm:char-span}, Theorem~\ref{thm:char-basis}).
